\documentclass[conference]{IEEEtran}
\IEEEoverridecommandlockouts

\usepackage{cite}
\usepackage{amsmath,amssymb,amsfonts}
\usepackage{algorithmic}
\usepackage{algorithm}
\usepackage{graphicx}
\usepackage{textcomp}
\usepackage{xcolor}
\usepackage{booktabs}
\usepackage{float}

\def\BibTeX{{\rm B\kern-.05em{\sc i\kern-.025em b}\kern-.08em
    T\kern-.1667em\lower.7ex\hbox{E}\kern-.125emX}}

\begin{document}

% ──────────────────────────────────────────────────────────
\title{DeadlockGNN: Real-Time Operating System Deadlock\\
Prediction Using Relational Graph Neural Networks\\
with Temporal Forecasting}
% ──────────────────────────────────────────────────────────

\author{
  \IEEEauthorblockN{Rishit Kamboj}
  \IEEEauthorblockA{Roll No.\ 2410110598\\rk846@snu.edu.in}
  \and
  \IEEEauthorblockN{Aadarsh Kumar}
  \IEEEauthorblockA{Roll No.\ 2410110388\\[FILL EMAIL]}
  \and
  \IEEEauthorblockN{Arjun Kapoor}
  \IEEEauthorblockA{Roll No.\ 2410110406\\[FILL EMAIL]}
  \and
  \IEEEauthorblockN{Shiven Agarwal}
  \IEEEauthorblockA{Roll No.\ 2410110321\\[FILL EMAIL]}
}

\maketitle

% ──────────────────────────────────────────────────────────
\begin{abstract}
Deadlock detection and prevention remain critical challenges in modern operating
systems and distributed computing environments. Classical approaches—including
Banker's Algorithm, probe-based detection, and Wait-For Graph (WFG) analysis—are
inherently reactive, can confirm a deadlock only after it has materialised, and
require complete a priori knowledge of resource needs. In this paper we present
\textbf{DeadlockGNN}, an end-to-end AI framework that reformulates deadlock
prediction as a graph classification problem over typed Resource Allocation
Graphs (RAGs). We apply a Relational Graph Convolutional Network (RGCN) for
instantaneous static deadlock classification and extend it with a Gated Recurrent
Unit (GRU) that processes sequences of eight consecutive RAG snapshots to
\emph{forecast} whether the next OS tick will be deadlocked. The static RGCN is
trained on 200,000 synthetic graphs; its learned weights are transferred into
the temporal encoder via cross-domain transfer learning. A live integration with
the xv6-riscv kernel running inside QEMU closes the loop between trained AI and
a real running operating system. Our temporal model achieves \textbf{93.74\%}
accuracy, \textbf{94.38\%} F1-score, and \textbf{0.9831} AUC-ROC at only
\textbf{2.9\,ms} per-sequence inference latency, demonstrating that GNN-based
proactive deadlock prediction is both feasible and computationally practical.
\end{abstract}

\begin{IEEEkeywords}
deadlock prediction, graph neural network, relational GCN, temporal forecasting,
operating systems, resource allocation graph, xv6, transfer learning
\end{IEEEkeywords}

% ──────────────────────────────────────────────────────────
\section{Introduction}
% ──────────────────────────────────────────────────────────

A deadlock occurs when a set of processes each holds at least one resource while
waiting for a resource held by another process in the set, forming a circular
dependency chain from which no process can escape without external
intervention~\cite{silberschatz2018}. The four necessary conditions—mutual
exclusion, hold-and-wait, no pre-emption, and circular wait—were formally
established by Coffman et al.~\cite{coffman1971} and remain the theoretical
foundation of all subsequent work.

Despite decades of research, deadlocks continue to be a practical concern in
concurrent operating systems, database engines, and large-scale distributed
systems. Traditional solutions fall into four families: (i)~\emph{prevention}
via strict resource-ordering protocols; (ii)~\emph{avoidance} via Banker's
Algorithm~\cite{dijkstra1965}; (iii)~\emph{detection} via cycle detection on
the Wait-For Graph; and (iv)~\emph{recovery} via process termination or
resource pre-emption. All four impose either significant runtime overhead or
require complete advance knowledge of resource demands—an assumption rarely
satisfied in real workloads.

Machine learning offers a fundamentally different paradigm: learn the
\emph{structural pre-conditions} of deadlock from data, without requiring
symbolic resource knowledge. Our literature survey~\cite{survey2024} identified
that recent work on Graph Neural Networks (GNNs) applied to system dependency
graphs achieves promising results. In particular, the DRIP (Dynamic Real-time
Integrated Predictive) framework~\cite{drip2024} demonstrated 99.9\% detection
accuracy by combining static program analysis with predictive modelling—however,
it operates on pre-defined rule sets rather than learned graph embeddings. GNN
approaches on resource allocation bipartite graphs~\cite{wang2021} show that
learned message-passing outperforms traditional greedy algorithms in dynamic
high-concurrency environments.

We build on these insights and make the following \textbf{contributions}:
\begin{enumerate}
  \item A typed 7-dimensional node feature encoding for RAGs that captures
        process role, contention pressure, and resource utilisation.
  \item A 3-layer \textbf{Relational GCN} with relation-specific weight matrices
        trained on 200,000 synthetic snapshots for $<$1\,ms static classification.
  \item A \textbf{Temporal RGCN-GRU} that forecasts deadlock at the next OS tick
        from 8 historical snapshots, using transfer learning from the static model.
  \item Live instrumentation of the \textbf{xv6-riscv} kernel with a real-time
        inference pipeline at 2.9\,ms latency.
  \item Monte-Carlo \textbf{Shapley attribution} for node-level explainability.
\end{enumerate}

% ──────────────────────────────────────────────────────────
\section{Literature Survey}
% ──────────────────────────────────────────────────────────

\subsection{Classical Deadlock Theory}

Coffman et al.~\cite{coffman1971} formalised the four necessary conditions for
deadlock and introduced the Wait-For Graph (WFG) as the canonical detection
structure. Dijkstra's Banker's Algorithm~\cite{dijkstra1965} provides a
polynomial-time safe-state check but requires complete knowledge of maximum
resource claims. Both approaches are reactive—they act on states that have
already formed, not on impending conditions.

\subsection{Predictive Deadlock Detection}

Mohamed et al.~\cite{drip2024} proposed the DRIP framework, a Dynamic
Real-time Integrated Predictive system that combines static analysis with
runtime monitoring to achieve 99.9\% detection accuracy in distributed settings.
While highly accurate, DRIP relies on hand-crafted rule integration rather than
learned representations, limiting its generalisability to unseen workload
patterns.

Gao et al.~\cite{gao2025} addressed deadlock detection specifically in deep
learning job schedulers using communication flow graphs, demonstrating that
graph-structured representations capture distributed deadlock conditions that
flat trace logs miss entirely.

\subsection{GNN-Based Resource Allocation and Scheduling}

Wang and Melchior~\cite{wang2021} modelled resource allocation as bipartite
graphs (tasks $\times$ resources) and showed that GNN-based allocation
strategies outperform linear programming and gradient-based baselines in
high-concurrency scenarios. This confirmed that GNNs can capture structural
resource contention in a way classical algorithms cannot.

Rout et al.~\cite{rout2021} extended distributed deadlock detection to cloud
computing environments, highlighting that probe-based algorithms suffer
prohibitive message overhead at scale—a gap that machine learning approaches
can address.

\subsection{Graph Neural Networks for Classification}

Kipf and Welling~\cite{kipf2017} established Graph Convolutional Networks
(GCNs) as the standard for semi-supervised node classification. Schlichtkrull
et al.~\cite{schlichtkrull2018} extended GCNs to relational (multi-edge-type)
graphs via R-GCN—the backbone of our model—enabling different weight matrices
per relation type. This is critical in our RAG setting where ``request'' and
``assignment'' edges carry opposite semantic meaning.

\subsection{Temporal Graph Learning}

Seo et al.~\cite{seo2018} combined GCNs with GRUs for spatio-temporal
forecasting, showing that sequential graph embeddings fed into a recurrent
model capture dynamics that single-snapshot GNNs cannot. Our temporal model
directly applies this insight to OS state evolution.

\subsection{Research Gap}

Existing work either (i) detects deadlocks reactively using WFG algorithms, or
(ii) applies ML to restricted domains (distributed schedulers, static code).
No prior system simultaneously: learns from raw OS RAG topology, forecasts
deadlock one tick ahead, and validates on a \emph{live running kernel}. DeadlockGNN
addresses all three.

% ──────────────────────────────────────────────────────────
\section{Proposed Methodology}
% ──────────────────────────────────────────────────────────

\subsection{Problem Formulation}

Let $G_t = (V_t, E_t)$ be the RAG at tick $t$, where $V_t = \mathcal{P}_t
\cup \mathcal{R}_t$ (processes and resources) and each edge carries a type
$\tau_e \in \{\textsc{req}, \textsc{asgn}\}$.

\noindent\textbf{Static task:} Given $G_t$, predict $y_t \in \{0,1\}$.

\noindent\textbf{Temporal task:} Given $\mathcal{S}_t = (G_{t-7}, \ldots, G_t)$,
predict $y_{t+1}$.

\subsection{Ground Truth: Wait-For Graph Cycle Detection}

\begin{algorithm}[H]
\caption{Deadlock Detection via WFG}
\begin{algorithmic}[1]
\REQUIRE RAG $G$
\STATE $\text{WFG} \leftarrow \emptyset$
\FOR{each resource $r$}
  \STATE $H \leftarrow \{p : (r \to p) \in G\}$;\quad
         $W \leftarrow \{p : (p \to r) \in G\}$
  \FOR{$w \in W$, $h \in H$, $w \neq h$}
    \STATE $\text{WFG} \mathrel{+}= (w \to h)$
  \ENDFOR
\ENDFOR
\STATE DFS on WFG; return \textit{true} iff back-edge found
\end{algorithmic}
\end{algorithm}

\subsection{Node Feature Encoding}

Each node $v$ is encoded as $\mathbf{f}_v \in \mathbb{R}^7$:
\begin{equation}
\mathbf{f}_v = \bigl[
  \mathbb{1}_P,\;
  \mathbb{1}_R,\;
  \tfrac{\deg v}{\max\!\deg},\;
  \tfrac{d^{in}}{d^{out}+1},\;
  |\mathcal{E}^{\text{req}}_v|,\;
  |\mathcal{E}^{\text{asgn}}_v|,\;
  \tfrac{|\text{holders}(v)|}{\text{cap}(v)}
\bigr]
\end{equation}

\subsection{Relational GCN (Static Model)}

RGCN propagation with $r \in \{\textsc{req}, \textsc{asgn}\}$:
\begin{equation}
\mathbf{h}_i^{(l+1)} = \sigma\!\left(
  \sum_r \sum_{j\in\mathcal{N}_r(i)}
  \frac{\mathbf{W}_r^{(l)}\,\mathbf{h}_j^{(l)}}{|\mathcal{N}_r(i)|}
\right)
\end{equation}
Three such layers are followed by global-add pooling and a two-layer MLP
classifier. The separate weight matrix $\mathbf{W}_r$ per relation is the
critical design choice: it lets the model learn that a process \emph{waiting
for} a resource has a structurally different role in deadlock formation than
one \emph{holding} it.

\begin{algorithm}[H]
\caption{RGCN Static Forward Pass}
\begin{algorithmic}[1]
\REQUIRE $\mathbf{X}\!\in\!\mathbb{R}^{n\times7}$, edges $E$, types $T$, batch $b$
\STATE $\mathbf{H} \leftarrow \text{ReLU}(\text{RGCNConv}_1(\mathbf{X},E,T))$
\STATE $\mathbf{H} \leftarrow \text{ReLU}(\text{RGCNConv}_2(\mathbf{H},E,T))$
\STATE $\mathbf{H} \leftarrow \text{RGCNConv}_3(\mathbf{H},E,T)$
\STATE $\mathbf{z} \leftarrow \text{GlobalAddPool}(\mathbf{H},b)$
\RETURN $\sigma(\text{MLP}(\mathbf{z}))$
\end{algorithmic}
\end{algorithm}

\subsection{Temporal RGCN-GRU (Predictive Model)}

\begin{algorithm}[H]
\caption{Temporal Forward Pass with Transfer Learning}
\begin{algorithmic}[1]
\REQUIRE Sequence $[G_0, \ldots, G_7]$
\STATE Load $\{$conv$_1$, conv$_2$, conv$_3\}$ from static checkpoint \COMMENT{Transfer}
\FOR{$t = 0$ \TO $7$}
  \STATE $\mathbf{e}_t \leftarrow \text{RGCN\_Encode}(G_t)$ \quad\COMMENT{shared weights, frozen}
\ENDFOR
\STATE $\mathbf{Z} \leftarrow [\mathbf{e}_0;\ldots;\mathbf{e}_7]\in\mathbb{R}^{8\times64}$
\STATE $\mathbf{h}_T \leftarrow \text{GRU}(\mathbf{Z},\;\text{layers}=2,\;\text{hidden}=128)$
\RETURN $\sigma(\text{MLP}(\mathbf{h}_T[-1]))$ \COMMENT{$P(\text{deadlock at }t{=}8)$}
\end{algorithmic}
\end{algorithm}

The RGCN encoder inherits spatial knowledge from 200,000 static training graphs
(transfer learning), so the GRU only needs to learn temporal dynamics from the
50,000-sequence dataset—drastically reducing convergence time.

\subsection{Dataset Generation}

Synthetic RAGs are generated by a parameterised OS simulator with:
Poisson process arrivals ($\lambda = 0.4$), configurable resource capacities
(up to 5~cores), hold durations of 3–10 ticks, and burst spawning (5\%
probability). Ground-truth deadlock labels are computed by Algorithm~1 on every
snapshot. Class balance (50/50) is enforced by rejection sampling.

\subsection{Live Kernel Integration}

The xv6-riscv kernel is instrumented in \texttt{spinlock.c} to emit:
\begin{verbatim}
[GNN_TRACE] LOCK_ACQUIRE P5 virtio_disk 1642
[GNN_TRACE] PROCESS_SLEEP P3 proc 1645
\end{verbatim}
A Python bridge reads events from a UNIX FIFO pipe and updates the live RAG
in $O(1)$ per event. Both models run inference every second on the current
kernel state.

% ──────────────────────────────────────────────────────────
\section{Results}
% ──────────────────────────────────────────────────────────

\subsection{Quantitative Evaluation}

\begin{table}[H]
\centering
\caption{Model Performance on Held-Out Test Sets}
\begin{tabular}{lccccc}
\toprule
\textbf{Model} & \textbf{Acc.} & \textbf{Prec.} & \textbf{Rec.} & \textbf{F1} & \textbf{AUC} \\
\midrule
RGCN (3k)          & 91.2\% & 91.8\% & 90.5\% & 91.1\% & 0.961 \\
RGCN (200k)        & 92.6\% & 93.1\% & 92.0\% & 92.5\% & 0.975 \\
\textbf{RGCN-GRU}  & \textbf{93.74\%} & \textbf{96.30\%} & \textbf{92.53\%} & \textbf{94.38\%} & \textbf{0.983} \\
\bottomrule
\end{tabular}
\end{table}

The temporal model evaluated over 10,000 test sequences produces the confusion
matrix:
$\bigl[\begin{smallmatrix} 4121 & 202 \\ 424 & 5253 \end{smallmatrix}\bigr]$
(TN, FP / FN, TP).
High precision (96.30\%) means very few false alarms—critical for a production
monitoring tool where unnecessary process termination is costly.

\subsection{Key Observations}

\textbf{RGCN vs.\ plain GCN:} Relation-aware message passing improves F1 by
$\sim$4\% over a GCN baseline, confirming that edge-type semantics are essential.

\textbf{Transfer learning:} Initialising the temporal encoder with static weights
reduces convergence from $>$100 epochs to 30 epochs and improves final F1 by 3.8\%.

\textbf{Latency:} At 2.9\,ms per-sequence inference, the system comfortably fits
within a 10\,ms OS scheduling quantum.

\textbf{Comparison with DRIP~\cite{drip2024}:} DRIP achieves 99.9\% on its own
benchmark using hand-crafted rules. DeadlockGNN, by contrast, learns purely
from graph topology and generalises to unseen process/resource configurations
without rule engineering.

% ──────────────────────────────────────────────────────────
\section{Conclusion}
% ──────────────────────────────────────────────────────────

DeadlockGNN demonstrates that OS deadlock prediction can be recast as a learned
graph classification problem with competitive accuracy and real-time latency.
By modelling OS state as typed Resource Allocation Graphs and applying Relational
GCN with semantically distinct edge-type weights, we achieve 91–93\% accuracy on
static snapshots. The temporal RGCN-GRU, which forecasts deadlock one tick ahead
using a sequence of eight historical RAG snapshots, achieves 93.74\% accuracy and
0.9831 AUC at 2.9\,ms—the first system to demonstrate both live kernel integration
and proactive forecasting in a unified framework.

Future directions include: (i)~extending to Linux via eBPF probes; (ii)~online
continual learning to adapt to shifting workloads; (iii)~graph-level attention
for interpretable cycle highlighting; and (iv)~multi-GPU distributed deadlock
prediction using heterogeneous resource graphs.

% ──────────────────────────────────────────────────────────
\begin{thebibliography}{00}

\bibitem{silberschatz2018}
A.\ Silberschatz, P.\ B.\ Galvin, and G.\ Gagne,
\textit{Operating System Concepts}, 10th ed.\ Hoboken, NJ: Wiley, 2018.

\bibitem{coffman1971}
E.\ G.\ Coffman, M.\ J.\ Elphick, and A.\ Shoshani,
``System deadlocks,''
\textit{ACM Computing Surveys}, vol.\ 3, no.\ 2, pp.\ 67--78, 1971.

\bibitem{dijkstra1965}
E.\ W.\ Dijkstra,
``Cooperating sequential processes,''
Tech.\ Rep.\ EWD-123, Eindhoven Univ.\ of Technology, 1965.

\bibitem{kipf2017}
T.\ N.\ Kipf and M.\ Welling,
``Semi-supervised classification with graph convolutional networks,''
in \textit{Proc.\ ICLR}, 2017.

\bibitem{schlichtkrull2018}
M.\ Schlichtkrull \textit{et al.},
``Modeling relational data with graph convolutional networks,''
in \textit{Proc.\ ESWC}, 2018, pp.\ 593--607.

\bibitem{seo2018}
Y.\ Seo, M.\ Defferrard, P.\ Vandergheynst, and X.\ Bresson,
``Structured sequence modeling with graph convolutional recurrent networks,''
in \textit{Proc.\ ICONIP}, 2018, pp.\ 362--373.

\bibitem{drip2024}
T.\ Mohamed, S.\ N.\ Kaif, V.\ M.\ Yuvaraj, and Y.\ Singh,
``Dynamic real-time deadlock detection framework using predictive analysis
and static integration,''
\textit{Int.\ J.\ Comput.\ Sci.\ Eng.}, 2024.

\bibitem{wang2021}
J.\ Wang and P.\ Melchior,
``Graph neural network resource allocation strategies,''
\textit{Astron.\ Comput.}, vol.\ 37, 2021.

\bibitem{rout2021}
S.\ K.\ Rout, B.\ Adhikari, and H.\ S.\ Behera,
``Distributed deadlock detection in cloud computing systems,''
\textit{J.\ Parallel Distrib.\ Comput.}, 2021.

\bibitem{gao2025}
Y.\ Gao \textit{et al.},
``Detecting communication deadlocks in deep learning jobs via graph analysis,''
in \textit{Proc.\ ICSE}, 2025.

\bibitem{survey2024}
R.\ Kamboj, A.\ Kumar, A.\ Kapoor, and S.\ Agarwal,
``Literature survey: Intelligent deadlock detection and resource allocation
using graph neural networks,''
Shiv Nadar Univ., Tech.\ Report, 2024.

\end{thebibliography}

\end{document}
